%%%%%%%%%%%%%%%%%%%%%%%%%%
%   DRAFT - Kevolution
%%%%%%%%%%%%%%%%%%%%%%%%%%

\documentclass[a4paper,11pt]{article}
\pdfoutput=1
\usepackage{jcappub}
\usepackage{bbold}
\usepackage[x11names]{xcolor}
\usepackage{colortbl,hhline}
\usepackage[english]{babel}
\usepackage[utf8]{inputenc}
\usepackage{amsmath}
\usepackage{color}
\usepackage{amsfonts}
\usepackage{graphicx}
\usepackage{amssymb}
\usepackage{eufrak}
\usepackage{etoolbox}
\usepackage{amsmath}
\usepackage{empheq}
\usepackage{cancel}
\usepackage[most]{tcolorbox}
\usepackage{float}              % Activate [H] option to place figure HERE
\usepackage{subcaption}
\usepackage[normalem]{ulem}
\usepackage{xspace}
\usepackage{adjustbox}
\newtcbox{\mymath}[1][]{%
    nobeforeafter, math upper, tcbox raise base,
    enhanced, colframe=yellow!30!black,
    colback=yellow!30, boxrule=1pt,
    #1}
%%%%%%%%%%%%%%%%%%%%%%%%%%%%

\def\be{\begin{equation}}
\def\ee{\end{equation}}
\def\bea{\begin{eqnarray}}
\def\eea{\end{eqnarray}}
\def\bean{\begin{eqnarray*}}
\def\eean{\end{eqnarray*}}
\def\cd{\cdot}
\def\vp{\varphi}
\def\l {\langle}
\def\re {\rangle}
\def \dd {\partial}
\def \ra {\rightarrow}
\def \la {\lambda}
\def \La {\Lambda}
\def \De {\Delta}
\def \DH {\Delta_{\rm HI}}
\newcommand{\de}{\delta}
\def \b {\beta}
\def \al {\alpha}
\def \ka {\kappa}
\def \Ga {\Gamma}
\def \ga {\gamma}
\def \si {\sigma}
\def \Si {\Sigma}
\def \ep {\epsilon}
\def \om {\omega}
\def \Om {\Omega}
\def \lap {\triangle}
\def \ep {\epsilon}


%%%%%%%%%%%%%%%%%%%%%%%%%%%%%%%%%%%
%Special definitions for this paper
%%%%%%%%%%%%%%%%%%%%%%%%%%%%%%%%%%%

\newcommand{\MyRed}{\color [rgb]{0.8,0,0}}
\newcommand{\MyGreen}{\color [rgb]{0,0.7,0}}
\newcommand{\MyBlue}{\color [rgb]{0,0,0.8}}
\newcommand{\MyBrown}{\color [rgb]{0.8,0.4,0.1}}
\newcommand{\MyPurple}{\color [rgb]{0.6,0.0,0.6}}
\def\GV#1{{\MyRed [GV: #1]}}
\def\RD#1{{\MyGreen [RD:  {\tt #1}]}} 
\def\RDt#1{{\MyGreen #1}}   
\def\mkt#1{{\MyBlue #1}}  
\def\mkr#1{{\MyBlue [MK: #1]}}  
\def\GF#1{{\MyPurple [GF: #1]}} 
\newcommand{\jat}[1]{{\color[rgb]{0,0.6,0}#1}}
\newcommand{\jar}[1]{{\color[rgb]{0,0.7,0} {\sf[#1]}}}
\newcommand{\FV}[1]{{\color[rgb]{0.8,0.33,0.13}#1}}

\newcommand{\ie}{\emph{i.e.}\xspace}
\newcommand{\cf}{\emph{cf.}\xspace}
\newcommand{\etal}{\emph{et al.}\xspace}
\newcommand{\eg}{\emph{e.g.}\xspace}

\newcommand{\Scal}{\mathcal S}
\newcommand{\DD}{\mathcal D}
\newcommand{\EE}{\mathcal E}
\newcommand{\MM}{\mathcal M}
\newcommand{\HH}{\mathcal H}

\newcommand{\Real}{\mathbb{R}}
\newcommand{\bn}{\boldsymbol{n}}
\newcommand{\bv}{\boldsymbol{v}}
\newcommand{\bx}{\boldsymbol{x}}
\newcommand{\bnabla}{\boldsymbol{\nabla}}
\newcommand{\bell}{\boldsymbol{\ell}}
\newcommand{\bal}{\boldsymbol{\alpha}}

\definecolor{bittersweet}{rgb}{1.0, 0.44, 0.37}
\newcommand{\fhc}[1]{{\textcolor{bittersweet}{\sf[#1]}}}
\newcommand{\fh}[1]{\textcolor{purple}{#1}}
\newcommand{\edit}[2]{\textcolor{red}{\sout{#1} {#2}}}  

\setlength{\arrayrulewidth}{.6mm}
\setlength{\tabcolsep}{7pt}
\renewcommand{\arraystretch}{1.8}

%Filippos commands
\newcommand{\MP}{M_{\rm Pl}}
\newcommand{\zetaf}{\zeta}

%Farbod
\newcommand{\lat}{{\mathrm i,\mathrm j,\mathrm k}}
\newcommand {\mf}{\mathrm}
\newcommand {\sn}{\mathrm n}
\newcommand {\gev}{\textit{gevolution}\xspace}
\newcommand {\kev}{{$k$-evolution}\xspace}
\newcommand{\class}{\texttt{CLASS}\xspace}
\newcommand{\hiclass}{\texttt{hi\_class}\xspace}

% for orcid
\usepackage{scalerel,tikz}
\usetikzlibrary{svg.path}
\definecolor{orcidlogocol}{HTML}{A6CE39}
\tikzset{orcidlogo/.pic={
 \fill[orcidlogocol] svg{M256,128c0,70.7-57.3,128-128,128C57.3,256,0,198.7,0,128C0,57.3,57.3,0,128,0C198.7,0,256,57.3,256,128z};
 \fill[white] svg{M86.3,186.2H70.9V79.1h15.4v48.4V186.2z}
 svg{M108.9,79.1h41.6c39.6,0,57,28.3,57,53.6c0,27.5-21.5,53.6-56.8,53.6h-41.8V79.1z M124.3,172.4h24.5c34.9,0,42.9-26.5,42.9-39.7c0-21.5-13.7-39.7-43.7-39.7h-23.7V172.4z}
 svg{M88.7,56.8c0,5.5-4.5,10.1-10.1,10.1c-5.6,0-10.1-4.6-10.1-10.1c0-5.6,4.5-10.1,10.1-10.1C84.2,46.7,88.7,51.3,88.7,56.8z};
}}
\newcommand\orcidicon[1]{\href{https://orcid.org/#1}{\mbox{\scalerel*{
\begin{tikzpicture}[yscale=-1,transform shape]
\pic{orcidlogo};
\end{tikzpicture}
}{|}}}}

%%%%%%%%%%%%%%%%%%%%%%%%%%%%%%%%%%%%%%%%%%
\title{$k$-essence}
\abstract{
}
\graphicspath{Images/}
\begin{document}
\maketitle
%%%%%%%%%%%%%%%%%
%%%INTRODUCTION
%%%%%%%%%%%%%%%%%

\section{Introduction}


\section{Action}

Starting from the covariant action of $k$-essence,
\be
S_{\mathrm{DE}}=\int d^{4} x \sqrt{-g} P(X, \phi)
\ee
Where the scalar field could be written as $\phi(t, \vec{x})=\bar{\phi}(t)+\delta \phi(t, \vec{x})$. Now we choose the time coordinate  such that $\delta \phi = 0$ and thus $t$ becomes a function of $\phi$. As a result, we have
\begin{align}
& X = - g^{\mu \nu} \partial_{\mu} \phi \partial_{\nu} \phi \to -g^{00} \dot{\bar{\phi}}(t)^2 \\ \nonumber &
\phi \to \bar{\phi}(t)
\end{align}
So the action is written as,
\be
S_{\mathrm{DE}}=\int d^{4} x \sqrt{-g} P(-g^{00} \dot{\bar{\phi}}(t)^2, \bar{\phi}(t))
\ee
where $g^{00} = -1 + \delta g^{00} $. So we can write the Lagrangian as,
\begin{align}
P \Big( \dot{\bar{\phi}}(t)^2 (1 - \delta g^{00}), \bar{\phi}(t) \Big)  & = P( \bar{X}, \bar{\phi}(t)) -  \bar{X} \frac{\partial P} {\partial \bar{X}}\bigg|_{X=\bar{X}}  \delta g^{00}  + \frac{ \bar{X}^2}{2!}\frac{\partial^2 P} {\partial \bar{X}^2}\bigg|_{X=\bar{X}}  ( \delta g^{00} )^2 \nonumber  \\ &  -  \frac{ \bar{X}^3}{3!}\frac{\partial^3 P} {\partial \bar{X}^3}\bigg|_{X=\bar{X}}  ( \delta g^{00} )^3 + \cdot \cdot \cdot
\end{align}
where we have defined $\bar{X} =\dot{\bar{\phi}}(t)^2 $. As a result the action could be written in the following form:
\be
S=\int d^{4} x \sqrt{-g}\left[\frac{M_{\mathrm{Pl}}^{2}}{2} R-\Lambda(t)-c(t) g^{00}+\frac{M_{2}^{4}(t)}{2}\left(\delta g^{00}\right)^{2}  +\frac{M_{3}^{4}(t)}{3!}\left(\delta g^{00}\right)^{3} + \cdot \cdot \cdot  \right].
\ee

%We want to add the $k$-essence to the EFT action, i.e.~we want to include terms linear in $\delta K$, which is at least of order $\epsilon^0$. Since we are going to keep in the action all terms of order $\epsilon^2$ and since  $\delta g^{00}$ is at least of order $\epsilon$, then we need only to include $\delta K \delta g^{00}$ and $\delta K (\delta g^{00})^2$ because $\delta K (\delta g^{00})^n$ with $n>2$ will be higher order than ${\cal O}(\epsilon^2)$. in conclusion, the action is
%\be
%\begin{split}
%\label{EFTaction}
%S=\int d^4 x \sqrt{-g} \bigg[ & \frac{\MP^2}{2} R -\Lambda (t) -c(t) g^{00} +\frac{M_2^4(t)}{2} \left (\delta g^{00}  \right )^2  \\
%  & - \frac{m_3^3(t)}{2}  \delta K \delta g^{00} - \frac{M_3^3(t)}{2} \delta K (\delta g^{00})^2 + ... \bigg] \;,
%\end{split}
%\ee
\subsection{Stueckelberg trick}
We want to make the scalar field perturbation explicit. We follow \cite{Hassani_2019} and perform  a diffeomorphism,  
\be
\label{time_diff}
t \to \tilde t(\eta) = t(\eta)+ \xi^0 (t (\eta) , \vec x) \;, \qquad \vec x \to \vec {\tilde x} = \vec x  \; ,
\ee
where

\be
\pi(t(\eta),\vec x) \doteq \xi^0 (t(\eta),\vec x) /a (t(\eta)) \;.
\ee 
Where $\eta$ is the conformal time,
\be
\eta:=\int \frac{d t}{a (t)}
\ee

\begin{equation}
g^{00} \rightarrow a^2 (1+ \HH \pi +  \pi')^2 g^{00}+2 a^2 (1+  \HH \pi +  \pi')  \partial_{\mu} \pi  g^{0 \mu}+ a(\eta)^2 g^{\mu \nu} \partial_{\mu} \pi \partial_{\nu} \pi
\end{equation}


We have to pay attention to expressing the cosmic time quantities in terms of conformal time when changing coordinates.
For instance,
\begin{align}
f(t) & \to \tilde f(\tilde t) = f(t(\eta) + a(\eta) \pi(\eta ,\vec x)) \\ \nonumber & 
 = f(\eta)+f^{\prime}(\eta) \pi(\eta, \vec{x})+\frac{1}{2} \Big(f^{\prime \prime}(\eta)- \HH(\eta) f^{\prime}(\eta)  \Big) \pi(\eta,\vec x)^{2} + ...
\end{align}


On the other hand, since the geometric quantities $g^{00}$ and $\delta K$ are constructed based on the metric in conformal time (based on our notation), so in order to perform the transformation on the vector $A^{\mu}$ we have
% Remember that although $\Lambda (t)$ quantities are expressed in terms of physical time. But our metric is defined (because of our definition) in terms of conformal time. That is why we have $g^{00}$ written in conformal time, so now we want to go from the conformal time to our new coordinate which is $\tilde t (\eta)$, i.e.,
% $\eta \to \eta + a(\eta) \pi (\eta, \vec x)$
\begin{equation}
A^{\mu} \rightarrow \tilde{A}^{\mu}=\frac{\partial  \tilde{x}^{\mu}}{\partial   x^{\nu}} A^{\nu}
\end{equation}
where results in,
\begin{align}
A^{0} \rightarrow \tilde{A}^{0}& =\frac{\partial  \tilde{x}^{0}}{\partial   x^{0}} A^{0}  + \frac{\partial  \tilde{x}^{0}}{\partial   x^{i}} A^{i} = \frac{\partial  (t(\eta) + a(\eta) \pi (\eta, \vec x))}{\partial   \eta} A^0 +  \frac{\partial  (t(\eta) + a(\eta) \pi (\eta, \vec x))}{\partial   x^i} A^i \\ \nonumber &=a(\eta) \big( 1 + \HH \pi(\eta,\vec x) + \pi'(\eta, \vec x) \Big) A^0 + a(\eta) \partial_i \pi(\eta, \vec x) A^i
\end{align}
In order to transform the vectors and convectors it is convenient to define the following matrices  \cite{Lewandowski:2016yce},
\begin{equation}
R_{\nu}^{\mu} \doteq \frac{\partial \tilde{x}^{\mu}}{\partial x^{\nu}}, \; \; (R^{-1})_{\nu}^{\mu} \doteq \frac{\partial \tilde{x}_{\mu}}{\partial x_{\nu}}
\end{equation}
where $R$ and $R^{-1}$ read as,
\begin{equation}
R_{\nu}^{\mu}=\left(\begin{array}{cc}
a(1+\HH \pi+\pi^{\prime}) & a \partial_{i} \pi \\
0 & \delta_{j}^{i}
\end{array}\right)_{\mu \nu}
\end{equation}
\begin{equation}
(R^{-1})_{\nu}^{\mu} =
\left(\begin{array}{cc}
  \frac{1}{a(1+\HH \pi+\pi^{\prime})} & - \frac{ \partial_{i} \pi}{1+ \HH \pi+\pi^{\prime}} \\
0 & \delta_{i}^{j}
\end{array}\right)_{\mu \nu}
\end{equation}
Since the scalar $A^{\mu} A_{\mu}$ should be invariant under the coordinate transformation,
\be
A^{\mu} A_{\mu} \rightarrow  (R^{-1})_{\mu}^{\nu} R_{\rho}^{\mu} A^{\rho} A_{\nu}
\ee
As expected we can easily verify that,
\be
(R^{-1})_{\mu}^{\nu} R_{\rho}^{\mu} = I
\ee
As a result in order to perform gauge transformation, for any upper and lower-index we need to apply $R$  and $R^{-1}$ respectively. So we have the following relations,
\begin{equation}
g^{\mu \nu} \rightarrow R_{\rho}^{\mu} R_{\sigma}^{\nu} g^{\rho \sigma}
\end{equation}
\begin{equation}
g_{\mu \nu} \rightarrow (R^{-1})^{\rho}_{\mu}(R^{-1})^{\sigma}_{\nu} g^{\rho \sigma}
\end{equation}
\begin{equation}
g^{00} \rightarrow a^2 (1+ \HH \pi +  \pi')^2 g^{00}+2 a^2 (1+  \HH \pi +  \pi')  \partial_{\mu} \pi  g^{0 \mu}+ a(\eta)^2 g^{\mu \nu} \partial_{\mu} \pi \partial_{\nu} \pi
\end{equation}
\begin{equation}
g^{0i} \rightarrow a(1+\HH \pi + \pi') g^{0i}+a(\eta) \partial_{j} \pi  g^{i j}
\end{equation}
\begin{equation}
g^{ij} \rightarrow g^{ij}
\end{equation}
\subsection{Action}
The cpvariant action after the Stueckelberg trick reads,
\be
\label{EFTaction_cov}
\begin{split}
S=\int d^4 x \sqrt{-g} \bigg\{ & \frac{\MP^2}{2} R -  \Lambda \left[ t(\eta)+ a(\eta) \pi \right] - c \left[ t(\eta)+ a(\eta) \pi \right]  a^2(\eta)  \big[ (1+ \HH \pi + \pi')^2 g^{00} \\
&+ 2 (1+\HH \pi + \pi') g^{0 \mu} \partial_\mu \pi + g^{\mu \nu} \partial_\mu \pi \partial_\nu \pi  \big]  +\frac{M_2^4 \left[ t(\eta)+ a(\eta) \pi \right]}{2}  \\
& \times a^4 (\eta)  \big[ (1+ \HH \pi + \pi')^2 g^{00} + 2 (1+\HH \pi + \pi') g^{0 \mu} \partial_\mu \pi + g^{\mu \nu} \partial_\mu \pi \partial_\nu \pi  - \bar g^{00}\big]^2 
\\ +\frac{M_3^4 \left[ t(\eta)+ a(\eta) \pi \right]}{3!}  
& \times a^6 (\eta)  \big[ (1+ \HH \pi + \pi')^2 g^{00} + 2 (1+\HH \pi + \pi') g^{0 \mu} \partial_\mu \pi + g^{\mu \nu} \partial_\mu \pi \partial_\nu \pi  - \bar g^{00}\big]^3      \bigg\} \;,
\end{split}
\ee
Here, we only focus on the last term in the EFT expansion, i.e.,
\be
\begin{split}
S=\int d^4 x \sqrt{-g} \bigg\{ & \frac{M_3^4 \left[ t(\eta)+ a(\eta) \pi \right]}{3!}  \\ & \times a^6 (\eta)  \big[ (1+ \HH \pi + \pi')^2 g^{00} + 2 (1+\HH \pi + \pi') g^{0 \mu} \partial_\mu \pi + g^{\mu \nu} \partial_\mu \pi \partial_\nu \pi  - \bar g^{00}\big]^3      \bigg\} \;,
\end{split}
\ee

%\begin{equation}
%g_{00} \rightarrow \frac{g_{00}}{a^2 (1+\HH \pi + \pi')^2} 
%\end{equation}
%\begin{equation}
%g_{0i} \rightarrow \frac{-\partial_i \pi g_{00}}{a(1+\HH \pi + \pi')^2}  +  \frac{ g_{0i}}{a(1+\HH \pi + \pi')}
%\end{equation}
%\begin{equation}
%g_{ij} \rightarrow  \frac{ \partial_{i} \pi \partial_{j} \pi}{(1+\HH \pi + \pi') ^2}\, g_{00} - \frac{a(\eta) \partial_{i} \pi  }{1+\HH \pi + \pi'} \, g_{j0}  +g_{ij} 
%\end{equation}
%The ADM quantities:
%
%\begin{equation}
%\begin{array}{c}
%g^{00}=-\frac{1}{N^{2}}, \quad g^{0 i}=\frac{N^{i}}{N^{2}}, \quad h^{i j}=g^{i j}+\frac{N^{i} N^{j}}{N^{2}} \\
%g_{00}=-N^{2}+h_{i j} N^{i} N^{j}, \quad N_{i}=h_{i j} N^{j}, \quad h_{i j}=g_{i j}
%\end{array}
%\end{equation}
%Here we only keep the relevant terms:
%\begin{equation}
%N^{i} \rightarrow N^{i}a (1-\HH{\pi} - \pi')+ a   N^{2} h^{i k} \partial_{k} \pi + \cdot \cdot \cdot
%\end{equation}
%\begin{equation}
%N \rightarrow N a \left(1- \HH{\pi} - \pi' + N^{i} \partial_{i} \pi+\frac{1}{2} a(\eta) N^{2} h^{i j} \partial_{i} \pi \partial_{i} \pi\right) + \cdot \cdot \cdot
%\end{equation}
%\begin{equation}
%h_{i j} \rightarrow h_{i j}-a(\eta) N_{i} \partial_{j} \pi- a(\eta) N_{j} \partial_{i} \pi- a(\eta)^2 N^{2} \partial_{i} \pi \partial_{j} \pi
%\end{equation}
%Now we are going to use the extrinsic curvature definition written in terms of the geometric objects on 3D hypersurfaces,
%

%
\bibliographystyle{JHEP}
\bibliography{Biblio}

\end{document}
